\chapter{Using VMEC}

\section{Representation of Flux Surface Geometry}
%
An inverse representation is used to describe the flux surface geometry in VMEC++.
This means that the (cylindrical) coordinates of the flux surface geometry
are expressed as functions of the flux coordinates.
\\
The three flux coordinates are:
\begin{itemize}
 \item $s$ is the normalized toroidal flux enclosed in a given flux surface.
       Its value is $0$ at the magnetic axis (in the core of the plasma) and $1$ at the last closed flux surface.
 \item $\theta$ is the poloidal angle-like coordinate.
       It increases by $2 \pi$ once around the torus the short way, i.e., wristband-like.
 \item $\varphi$ is the toroidal angle
       and it coincides with the angle coordinate in a cylindrical coordinate system.
       It increases by $2 \pi$ once around the torus the long way.
\end{itemize}
The full toroidal angle $\varphi$ is assumed to be made up of $n_\textrm{fp}$ identical, rotated, field periods of the plasma.
The toroidal coordinate in each field period is then
\begin{equation}
 \zeta = n_\textrm{fp} \varphi \, ,
\end{equation}
such that $\zeta$ goes from $0$ to $2 \pi$ over a single field period.

Having the flux coordinates $(s, \theta, \zeta)$, the cylindrical coordinates~$(R, \varphi, Z)$ of points on the flux surfaces
are given by:
\begin{align}
 R       =&\, R(s, \theta, \zeta) \nonumber \\
 \varphi =&\, \zeta / n_\textrm{fp} \\
 Z       =&\, Z(s, \theta, \zeta) \nonumber \, .
\end{align}
A set of discrete flux surfaces is considered in VMEC++.
This is reflected in the radial coordinate $s$, which only takes on discrete values in VMEC++.
The number of flux surfaces is an input parameter, which is called \texttt{ns}.
The individual flux surfaces are then (usually) located at equal increments in $s$:
\begin{equation}
 s_j = \frac{j}{\texttt{ns} - 1} \textrm{ for } j \in \{0, 1, ..., \texttt{ns} - 1\} \, .
\end{equation}
Toroidal symmetry is taken into account for the number of field periods, $n_\textrm{fp}$.

For each discrete flux surface, the coordinates $R$ and $Z$ on it are represented as two-dimensional Fourier series:
\begin{align}
 R(s_j, \theta, \zeta) =&\, \sum\limits_{m, n}
   \left[ \hat{R}_{j, m, n}^{\cos} \cos(m \theta - n \zeta) +
          \hat{R}_{j, m, n}^{\sin} \sin(m \theta - n \zeta) \right] \label{eqn:r_full} \\
 Z(s_j, \theta, \zeta) =&\, \sum\limits_{m, n}
   \left[ \hat{Z}_{j, m, n}^{\cos} \cos(m \theta - n \zeta) +
          \hat{Z}_{j, m, n}^{\sin} \sin(m \theta - n \zeta) \right] \label{eqn:z_full} \, .
\end{align}
In case of stellarator symmetry, half of the Fourier coefficients can be omitted
and the flux surface geometry is given by:
\begin{align}
 R(s_j, \theta, \zeta) =&\, \sum\limits_{m, n}
   \hat{R}_{j, m, n}^{\cos} \cos(m \theta - n \zeta) \label{eqn:r_symm} \\
 Z(s_j, \theta, \zeta) =&\, \sum\limits_{m, n}
   \hat{Z}_{j, m, n}^{\sin} \sin(m \theta - n \zeta) \textsl{}\label{eqn:z_symm} \, .
\end{align}
The summation over $m$ and $n$ goes as follows.
The poloidal mode number $m$ goes from $0$ to $\texttt{mpol} - 1$.
For $m = 0$, only non-negative toroidal mode numbers are taken into account,
thus $n \in \{0, 1, ..., \texttt{ntor}\}$.
For $m > 0$, all toroidal mode numbers are taken into account,
thus $n \in \{-\texttt{ntor}, ..., -1, 0, 1, ..., \texttt{ntor}\}$.
This scheme serves to remove the redundancy (and associated consistency requirement) among the Fourier coefficients with $m = 0$.

The geometry of the magnetic axis is represented in a similar functional form.
The axis is a line along the toroidal direction and thus has to variation in the poloidal direction.
Therefore, it is only parameterized in the toroidal direction.

The Fourier series for $R$ and $Z$ of the flux surface geometry
are typically implemented as follows:
\begin{align}
 R(s_j, \theta, \varphi) =& \sum_{\texttt{mn} = 0}^{\texttt{mnmax}}
   \left[   R_\texttt{mn}^\mathrm{cos}(s_j) \cos(\texttt{xm[mn]} \theta - \texttt{xn[mn]} \varphi)
          + R_\texttt{mn}^\mathrm{sin}(s_j) \sin(\texttt{xm[mn]} \theta - \texttt{xn[mn]} \varphi) \right] \\
 Z(s_j, \theta, \varphi) =& \sum_{\texttt{mn} = 0}^{\texttt{mnmax}}
   \left[   Z_\texttt{mn}^\mathrm{sin}(s_j) \sin(\texttt{xm[mn]} \theta - \texttt{xn[mn]} \varphi)
          + Z_\texttt{mn}^\mathrm{cos}(s_j) \cos(\texttt{xm[mn]} \theta - \texttt{xn[mn]} \varphi) \right]
\end{align}
on, for example, the $j$-th flux surface with flux surface label (i.e., radial coordinate)~$s_j$.
%
The dataset names of these Fourier coefficients in the \texttt{wout} group
of the main VMEC++ output file are listed in Tab.~\ref{tab:fc_rz}.
\begin{table}[h]
 \centering
 \begin{tabular}{|c|c|c|}
  \hline
  symbol & dataset name & stellarator-symmetric? \\
  \hline
  $R_\texttt{mn}^\mathrm{cos}$ & \texttt{rmnc} & y \\
  $Z_\texttt{mn}^\mathrm{sin}$ & \texttt{zmns} & y \\
  \hline
  $R_\texttt{mn}^\mathrm{sin}$ & \texttt{rmns} & n \\
  $Z_\texttt{mn}^\mathrm{cos}$ & \texttt{zmnc} & n \\
  \hline
 \end{tabular}
 \caption{Dataset names for the Fourier coefficient arrays in the \texttt{wout} group
          that describe the cylindrical coordinates~$R$ and $Z$ of the flux-surface geometry.}
 \label{tab:fc_rz}
\end{table}

\section{Representation of Magnetic Field}
The magnetic field~$\mathbf{B}$ in the VMEC++ outputs
is available in both a contravariant form
and a covariant form.
The contravariant form is defined as follows:
\begin{equation}
 \mathbf{B} = B^\theta \hat{\mathbf{e}}_\theta + B^\zeta \hat{\mathbf{e}}_\zeta
\end{equation}
and the covariant form is defined as follows:
\begin{equation}
 \mathbf{B} = B_s \nabla s + B_\theta \nabla \theta + B_\zeta \nabla \zeta \, .
\end{equation}
The components $B^\theta$, $B^\zeta$, $B_\theta$ and $B_\zeta$
are available on the radial half-grid.
The component $B_s$ is available on the radial full-grid.
%
All these components are expressed as two-dimensional Fourier series
to represent their scalar, real values on a discrete set of flux surfaces:
\begin{equation}
 X(s_j, \theta, \varphi) =
 \sum\limits_{m} \sum\limits_{n}
   \left[   \hat{X}_{mn}^\textrm{cos}(s_j) \cos(m \theta - n n_\textrm{fp} \varphi)
          + \hat{X}_{mn}^\textrm{sin}(s_j) \sin(m \theta - n n_\textrm{fp} \varphi) \right]
\end{equation}
where $X$ is one of $\{B^\theta, B^\zeta, B_s, B_\theta, B_\zeta\}$
and $\hat{X}_{mn}^\textrm{cos}$ ($\hat{X}_{mn}^\textrm{sin}$)
are the Fourier coefficients for the cosine (sine) basis functions, respectively,
on the $j$-th flux surface with flux surface label (i.e., radial coordinate)~$s_j$.
The number of toroidal field periods is denoted by~$n_\textrm{fp}$ here.
The poloidal angle-like coordinate~$\theta$ spans from $0$ to $2 \pi$
the short way around the torus.
The toroidal angle (indentical with the cylindrical angle)~$\varphi$
spans from $0$ to $2 \pi$ once around the machine.

The summation over $m$ and $n$ is to be taken over all relevant mode number combinates
needed to represent the magnetic field components.
Requiring that the components are real-valued
allows to make use of certain symmetries, reducing the number of required mode number
combinations to take into account.
%
In VMEC++, a linear mode index is used in conjunction with mode number arrays
for convenient one-dimensional summation:
\begin{align}
 X(s_j, \theta, \varphi) =
 \sum\limits_{\texttt{mn}=0}^{\texttt{mnmax\_nyq} - 1}
   \Bigl[ &\phantom{+}~ \hat{X}_\texttt{mn}^\textrm{cos}(s_j) \cos(\texttt{xm}[\texttt{mn}] \theta - \texttt{xn}[\texttt{mn}] \varphi) \nonumber \\
         ~&         +   \hat{X}_\texttt{mn}^\textrm{sin}(s_j) \sin(\texttt{xm}[\texttt{mn}] \theta - \texttt{xn}[\texttt{mn}] \varphi) \Bigr] \, .
\end{align}
%
It is noted that a larger number of Fourier coefficients~($\texttt{mnmax\_nyq}$)
is used to describe the magnetic field components,
than is used to describe the flux surface geometry~($\texttt{mnmax}$).
%
The dataset names of these Fourier coefficients in the \texttt{wout} group
of the main VMEC++ output file are listed in Tab.~\ref{tab:fc_gb}.
Note that the notation~$B = \vert\mathbf{B}\vert$
has been used for the magnetic field strength.
\begin{table}[h]
 \centering
 \begin{tabular}{|c|c|c|}
  \hline
  symbol & dataset name & stellarator-symmetric? \\
  \hline
  $\sqrt{g}_\texttt{mn}^\mathrm{cos}$   & \texttt{gmnc}     & y \\
  $B_\texttt{mn}^\mathrm{cos}$          & \texttt{bmnc}     & y \\
  $B_{\theta,\texttt{mn}}^\mathrm{cos}$ & \texttt{bsubumnc} & y \\
  $B_{\zeta,\texttt{mn}}^\mathrm{cos}$  & \texttt{bsubvmnc} & y \\
  $B_{s,\texttt{mn}}^\mathrm{sin}$      & \texttt{bsubsmns} & y \\
  $B_\texttt{mn}^{\theta,\mathrm{cos}}$ & \texttt{bsupumnc} & y \\
  $B_\texttt{mn}^{\zeta,\mathrm{cos}}$  & \texttt{bsupvmnc} & y \\
  \hline
  $\sqrt{g}_\texttt{mn}^\mathrm{sin}$   & \texttt{gmns}     & n \\
  $B_\texttt{mn}^\mathrm{sin}$          & \texttt{bmns}     & n \\
  $B_{\theta,\texttt{mn}}^\mathrm{sin}$ & \texttt{bsubumns} & n \\
  $B_{\zeta,\texttt{mn}}^\mathrm{sin}$  & \texttt{bsubvmns} & n \\
  $B_{s,\texttt{mn}}^\mathrm{cos}$      & \texttt{bsubsmnc} & n \\
  $B_\texttt{mn}^{\theta,\mathrm{sin}}$ & \texttt{bsupumns} & n \\
  $B_\texttt{mn}^{\zeta,\mathrm{sin}}$  & \texttt{bsupvmns} & n \\
  \hline
 \end{tabular}
 \caption{Dataset names for the Fourier coefficient arrays in the \texttt{wout} group
          that describe the Jacobian, the magnetic field strength,
          and the co- and contravariant magnetic field components.}
 \label{tab:fc_gb}
\end{table}


\section{Re-Computing the Rotational Transform Profile}
The contravariant magnetic field components are defined as follows in VMEC~\cite{lee_1988}:
\begin{align}
 B^{\theta}  =&\, \frac{\Phi'}{\sqrt{g}} \left( \iota - \frac{\partial \lambda}{\partial \varphi} \right) \label{eqn:bsupu} \\
 B^{\varphi} =&\, \frac{\Phi'}{\sqrt{g}} \left(    1  + \frac{\partial \lambda}{\partial \theta } \right) \label{eqn:bsupv} \, .
\end{align}
The covariant magnetic field components can be computed from these using the metric elements:
\begin{align}
 B_{\theta } =&\, g_{\theta  \theta} B^{\theta} + g_{\theta  \varphi} B^{\varphi} \label{eqn:bsubu} \\
 B_{\varphi} =&\, g_{\theta \varphi} B^{\theta} + g_{\varphi \varphi} B^{\varphi} \label{eqn:bsubv} \, .
\end{align}
The radial profile of enclosed net toroidal current is given by:
\begin{align}
 \frac{\mu_0}{2 \pi} I_\zeta
 =&\, \langle B_\theta \rangle \label{eqn:curtor} \, .
\end{align}
where $\langle . \rangle$ is an average over the corresponding flux surface:
\begin{equation}
 \langle B_\theta \rangle = \frac{1}{(2 \pi)^2} \int\limits_{0}^{2 \pi} \int\limits_{0}^{2 \pi}
        B_\theta(s, \theta, \varphi) \,\mathrm{d}\theta \,\mathrm{d}\varphi \, .
\end{equation}
%
Now, \eqn{bsupu} and \eqn{bsupv} are inserted into \eqn{bsubu}:
\begin{align}
 B_{\theta } =&\, \frac{\Phi'}{\sqrt{g}} \left[
       g_{\theta  \theta}  \left( \iota - \frac{\partial \lambda}{\partial \varphi} \right)
     + g_{\theta  \varphi} \left(    1  + \frac{\partial \lambda}{\partial \theta } \right) \right] \, .
\end{align}
This in turn is inserted into \eqn{curtor}:
\begin{equation}
 \frac{\mu_0}{2 \pi} I_\zeta
 = \left\langle
     \frac{\Phi'}{\sqrt{g}} \left[
       g_{\theta  \theta}  \left( \iota - \frac{\partial \lambda}{\partial \varphi} \right)
     + g_{\theta  \varphi} \left(    1  + \frac{\partial \lambda}{\partial \theta } \right) \right] \right\rangle \, .
\end{equation}
The linearity of the integral allows to split this term into several parts.
Also, the quantities $\Phi'$ and $\iota$ only depend on the radial coordinate
and can be pulled out of the surface-averaging operator.
It thus follows:
\begin{equation}
 \frac{\mu_0}{2 \pi} \frac{I_\zeta}{\Phi'}
 =
 \iota \left\langle \frac{g_{\theta  \theta}}{\sqrt{g}} \right\rangle
 + \left\langle - \frac{g_{\theta  \theta }}{\sqrt{g}} \frac{\partial \lambda}{\partial \varphi}
                + \frac{g_{\theta  \varphi}}{\sqrt{g}} \frac{\partial \lambda}{\partial \theta } \right\rangle
 + \left\langle \frac{g_{\theta  \varphi}}{\sqrt{g}} \right\rangle \, .
\end{equation}
This can be re-arranged for the iota profile:
\begin{equation}
 \iota \left\langle \frac{g_{\theta  \theta}}{\sqrt{g}} \right\rangle
 =
 \frac{\mu_0}{2 \pi} \frac{I_\zeta}{\Phi'}
 - \left\langle - \frac{g_{\theta  \theta }}{\sqrt{g}} \frac{\partial \lambda}{\partial \varphi}
                + \frac{g_{\theta  \varphi}}{\sqrt{g}} \frac{\partial \lambda}{\partial \theta } \right\rangle
 - \left\langle \frac{g_{\theta  \varphi}}{\sqrt{g}} \right\rangle \label{eqn:iota_eqn} \, .
\end{equation}
The quantity $\lambda$ is represented as a Fourier series in VMEC
(with only $\sin$-coefficients assuming stellarator-symmetry):
\begin{equation}
 \lambda(s, \theta, \varphi)
 =
 \sum\limits_{m,n}
   \hat{\lambda}^{\sin}_{mn}(s)
   \sin\left(m \theta - n n_\textrm{fp} \varphi \right) \, .
\end{equation}
This yields analytical derivatives in both tangential directions:
\begin{align}
 \frac{\partial \lambda}{\partial \theta}
 =&\, \sum\limits_{m,n}
        m
        \hat{\lambda}^{\sin}_{mn}
        \cos\left(m \theta - n n_\textrm{fp} \varphi \right) \label{eqn:lu} \\
 \frac{\partial \lambda}{\partial \varphi}
 =&\, \sum\limits_{m,n}
        (-n n_\textrm{fp})
        \hat{\lambda}^{\sin}_{mn}
        \cos\left(m \theta - n n_\textrm{fp} \varphi \right) \label{eqn:lv} \, .
\end{align}
\eqn{lu} and \eqn{lv} are now inserted into \eqn{iota_eqn}.
The Fourier coefficients can be regarded as functions
of only the radial coordinate $s$ and can therefore be pulled out of the surface averaging operator.
\begin{equation}
 \iota \left\langle \frac{g_{\theta  \theta}}{\sqrt{g}} \right\rangle
 =
 - \sum\limits_{m,n} \left\langle \left(- \frac{g_{\theta  \theta }}{\sqrt{g}} (-n n_\textrm{fp})
                                        + \frac{g_{\theta  \varphi}}{\sqrt{g}}   m
                                  \right) \cos\left(m \theta - n n_\textrm{fp} \varphi \right)
                      \right\rangle
                      \hat{\lambda}^{\sin}_{mn}
 - \left\langle \frac{g_{\theta  \varphi}}{\sqrt{g}} \right\rangle
 + \frac{\mu_0}{2 \pi} \frac{I_\zeta}{\Phi'} \label{eqn:full_iota} \, .
\end{equation}
A matrix $\mathbf{A}$ with entries $a_{mn}$ and vectors $b$ and $c$ are now introduced with
\begin{equation}
 a_{mn} = \left\langle \left(  \frac{g_{\theta  \theta }}{\sqrt{g}} n n_\textrm{fp}
                                   + \frac{g_{\theta  \varphi}}{\sqrt{g}} m
                             \right) \cos\left(m \theta - n n_\textrm{fp} \varphi \right) \right\rangle
          /
          \left\langle \frac{g_{\theta  \theta}}{\sqrt{g}} \right\rangle \label{eqn:a_mat}
\end{equation}
as well as
\begin{equation}
 b = \left\langle \frac{g_{\theta  \varphi}}{\sqrt{g}} \right\rangle
     /
     \left\langle \frac{g_{\theta  \theta}}{\sqrt{g}} \right\rangle \label{eqn:b_vec} \, .
\end{equation}
and
\begin{equation}
 c = \frac{\mu_0}{2 \pi} \frac{I_\zeta}{\Phi'}
     /
     \left\langle \frac{g_{\theta  \theta}}{\sqrt{g}} \right\rangle \label{eqn:c_vec} \, .
\end{equation}
This allows to write \eqn{full_iota} as a linear system of equations:
\begin{equation}
 \iota = -\sum\limits_{m,n} a_{mn} \hat{\lambda}^{\sin}_{mn} - b + c \label{eqn:iota_from_linear_eqns} \, .
\end{equation}
%

\subsection{Numerical Implementation}
The objective of this section is go detail out the computation of
the quantites from \eqn{a_mat}, \eqn{b_vec}, \eqn{c_vec} for evaluation of $\iota$
from \eqn{iota_from_linear_eqns}.
%
The rotational transform profile is on the half-grid in VMEC.
Also, the metric elements, the Jacobian and the enclosed toroidal current as well as the toroidal flux derivative
are available on the half-grid in VMEC.
The state variable $\lambda$ is on the full-grid in VMEC
and needs to be averaged onto the half-grid before being used in \eqn{iota_from_linear_eqns}.




\subsection{Comparison with Mercier's Expression}
Mercier's expression for the rotational transform on the magnetic axis
of a stellarator is as follows:
\begin{equation}
 \iota =
 \frac{1}{2 \pi} \int\limits_{0}^{L}
   \frac{1}{\cosh \eta} \left[
       \frac{\mu_0 J}{2 B_0}
     - \left(\cosh \eta - 1 \right) d'
     - \tau
  \right]  \,\mathrm{d}l - N
\end{equation}
The three contributions in this expression
are due to the toroidal current,
a poloidal rotation of the poloidal cross section of the flux surfaces close to the axis
and torsion of the magnetic axis.
%
$\lambda$ is identically zero on the magnetic axis,
because the magnetic field is purely toroidal there and has no variation over the surface
of the axis, because the axis has no surface.
%
It is expected that the expression in \eqn{iota_from_linear_eqns} should approach Mercier's result
when moving towards the magnetic axis.
%
Here, $\mu_0 J$ is the current density on the magnetic axis.




\section{Current Density from VMEC}
The current density from VMEC is computed using Ampere's law:
\begin{equation}
 \mathbf{j} = \frac{1}{\mu_0} \nabla \times \mathbf{B} \, .
\end{equation}
%
If the curl of a covariant representation of a field is computed,
this results in a contravariant representation of the resulting field and vice versa~\cite{dHaseleer}.
%
Therefore, the covariant magnetic field components $B_i$ are needed
to compute the contravariant current density components $j^i$.
%
The covariant magnetic field components~$(B_s, B_\theta, B_\varphi)$ represent the equilibrium magnetic field $\mathbf{B}$ as follows:
\begin{equation}
 \mathbf{B} = B_s \nabla s + B_\theta \nabla \theta + B_\zeta \nabla \zeta \, .
\end{equation}
%
They are available in the output of VMEC as two-dimensional Fourier series on the half-grid:
\begin{align}
 B_s      (s_{j+0.5}, \theta, \varphi) =& \sum_{m,n} \left[  \hat{B}_{s,      mn}^\mathrm{cos} (s_{j+0.5}) \,\mathrm{cos}(m \theta - n n_\mathrm{fp} \varphi)
                                                            + \hat{B}_{s,      mn}^\mathrm{sin} (s_{j+0.5}) \,\mathrm{sin}(m \theta - n n_\mathrm{fp} \varphi) \right] \nonumber \\
 B_\theta (s_{j+0.5}, \theta, \varphi) =& \sum_{m,n} \left[  \hat{B}_{\theta, mn}^\mathrm{cos} (s_{j+0.5}) \,\mathrm{cos}(m \theta - n n_\mathrm{fp} \varphi)
                                                            + \hat{B}_{\theta, mn}^\mathrm{sin} (s_{j+0.5}) \,\mathrm{sin}(m \theta - n n_\mathrm{fp} \varphi) \right] \\
 B_\varphi  (s_{j+0.5}, \theta, \varphi) =& \sum_{m,n} \left[  \hat{B}_{\varphi,  mn}^\mathrm{cos} (s_{j+0.5}) \,\mathrm{cos}(m \theta - n n_\mathrm{fp} \varphi)
                                                            + \hat{B}_{\varphi,  mn}^\mathrm{sin} (s_{j+0.5}) \,\mathrm{sin}(m \theta - n n_\mathrm{fp} \varphi) \right] \nonumber
\end{align}
with the $m,n$ summation as detailed in Sec.~\ref{sec:torcoords_numbering},
albeit more Fourier coefficients are retained in this expansion
in order to satisfy the Nyquist sampling theorem.
%
Note that, in contrast to the contravariant magnetic field representation,
the radial component $B_s \neq 0$ in the covariant representation.
%
Given the covariant magnetic field representation, the curl of it takes a particurarly simple form~\cite{dHaseleer}:
\begin{equation}
 \mu_0 \mathbf{j}^i = \frac{1}{\sqrt{g}} \sum_k \left(  \frac{\partial B_j}{\partial i}
                                             - \frac{\partial B_i}{\partial j} \right)
                                       \hat{\mathbf{e}}_k, \quad (i,j,k) ~ \mathrm{cyc.} ~ (s, \theta, \varphi)
\end{equation}
with $\sqrt{g}$ being the Jacobian of the coordinate transform between flux coordinates and cylindrical coordinates.
%
The current density is, in practise, often computed to act as a source term in the Biot-Savart volume integrals.
There, the Jacobian~$\sqrt{g}$ is required in the integral.
It is therefore customary to compute~$\mathbf{j} \sqrt{g}$:
\begin{equation}
 j^i \sqrt{g} = \frac{1}{\mu_0} \left( \frac{B_k}{\partial u^j} - \frac{B_j}{\partial u^k} \right) , \quad (i,j,k) \textrm{ cyc. } (s,\theta,\zeta) \, .
\end{equation}
%
The concrete contravariant components of the current density are thus:
\begin{align}
 j^s      \sqrt{g} =&\, \frac{1}{\mu_0} \left( \frac{\partial B_\zeta }{\partial \theta} - \frac{\partial B_\theta}{\partial \zeta } \right) \nonumber \\
 j^\theta \sqrt{g} =&\, \frac{1}{\mu_0} \left( \frac{\partial B_s     }{\partial \zeta } - \frac{\partial B_\zeta }{\partial s     } \right) \label{eqn:currx} \\
 j^\zeta  \sqrt{g} =&\, \frac{1}{\mu_0} \left( \frac{\partial B_\theta}{\partial s     } - \frac{\partial B_s     }{\partial \theta} \right) \nonumber
\end{align}
%
The current density vector $\mathbf{j}$ is constrained to be tangential to flux surfaces in VMEC++.
%
Therefore, the current density can only have non-zero components in the tangential $\theta$ and $\varphi$ directions and $j^s = 0$ is assumed.
%
This allows to formulate a magnetic differential equation:
\begin{align}
 \frac{\partial B_\zeta }{\partial \theta} - \frac{\partial B_\theta}{\partial \zeta } = 0 \, .
\end{align}
%
We only consider $j^\theta$ and $j^\zeta$ in the following.
%
The two-dimensional Fourier series representation of the covariant magnetic field components
allows to perform the tangential derivatives in \eqn{currx} analytically:
\begin{align}
 \frac{\partial B_i}{\partial \theta}(s, \theta, \zeta)
 =&\, \sum\limits_{m,n} \left[
   \hat{B}^\textrm{cos}_{i,m,n}(s) m (-\sin)(m \theta - n \zeta) +
   \hat{B}^\textrm{sin}_{i,m,n}(s) m   \cos (m \theta - n \zeta)
 \right] \nonumber \\ ~
 =&\, \sum\limits_{m,n} \left[
   \left(-m \hat{B}^\textrm{cos}_{i,m,n}(s)\right) \sin(m \theta - n \zeta) +
   \left( m \hat{B}^\textrm{sin}_{i,m,n}(s)\right) \cos(m \theta - n \zeta)
 \right] \\
 \frac{\partial B_i}{\partial \zeta}(s, \theta, \zeta)
 =&\, \sum\limits_{m,n} \left[
   \hat{B}^\textrm{cos}_{i,m,n}(s) (-n) (-\sin)(m \theta - n \zeta) +
   \hat{B}^\textrm{sin}_{i,m,n}(s) (-n) \cos(m \theta - n \zeta)
 \right] \nonumber \\ ~
 =&\, \sum\limits_{m,n} \left[
   \left( n \hat{B}^\textrm{cos}_{i,m,n}(s)\right) \sin(m \theta - n \zeta) +
   \left(-n \hat{B}^\textrm{sin}_{i,m,n}(s)\right) \cos(m \theta - n \zeta)
 \right] \, .
\end{align}
%
The quantities $j^i \sqrt{g} \equiv J^i$ are also Fourier-decomposed:
\begin{equation}
 J^i(s, \theta, \zeta) = \sum\limits_{m,n} \left[
   \hat{J}^{i,\textrm{cos}}_{m,n}(s) \cos(m \theta - n \zeta) +
   \hat{J}^{i,\textrm{sin}}_{m,n}(s) \sin(m \theta - n \zeta)
 \right] \, .
\end{equation}
%
The orthogonality of the Fourier basis now allows to express \eqn{currx}
for each of the coefficients independently:
\begin{align}
 \hat{J}^{\theta,\textrm{cos}}_{m,n}(s) =&\, \frac{1}{\mu_0} \left(
   -n \hat{B}^\textrm{sin}_{s,m,n}(s) - \frac{\partial \hat{B}^\textrm{cos}_{\zeta,m,n}}{\partial s}
 \right) \\
 \hat{J}^{\theta,\textrm{sin}}_{m,n}(s) =&\, \frac{1}{\mu_0} \left(
   ~ \,\, n \hat{B}^\textrm{cos}_{s,m,n}(s) - \frac{\partial \hat{B}^\textrm{sin}_{\zeta,m,n}}{\partial s}
 \right) \\
 \hat{J}^{\zeta,\textrm{cos}}_{m,n}(s) =&\, \frac{1}{\mu_0} \left(
   \frac{\partial \hat{B}^\textrm{cos}_{\theta,m,n}}{\partial s} - m \hat{B}^\textrm{sin}_{s,m,n}(s)
 \right) \\
 \hat{J}^{\zeta,\textrm{sin}}_{m,n}(s) =&\, \frac{1}{\mu_0} \left(
   \frac{\partial \hat{B}^\textrm{sin}_{\theta,m,n}}{\partial s} + m \hat{B}^\textrm{cos}_{s,m,n}(s)
 \right)
\end{align}
%
If the current density components are desired to be computed on the full-grid,
the radial derivatives can make use of the fact that $B_\theta$ and $B_\zeta$
are computed on the half-grid in VMEC++.
%
Derivatives with respect to $s$ are computed as finite differences across the discrete flux surfaces.
%
Regularization terms have to be included in order to
reduce the impact of catastrophic cancellation near the magnetic axis~\cite{hirshman_vmec_currents}.
%
Only the stellarator-symmetric current density components~$\hat{j}_{mn}^{\theta, \mathrm{cos}}$
and~$\hat{j}_{mn}^{\varphi,  \mathrm{cos}}$ are considered in the following.
%
The non-stellarator-symmetric components follow analogously.
%
The radial derivatives~$\partial \hat{B}_{\theta, mn}^\mathrm{cos} / \partial s$
and~$\partial \hat{B}_{\varphi, mn}^\mathrm{cos} / \partial s$ are computed as follows:
\begin{align}
\frac{\partial \hat{B}_{\theta, mn}^\mathrm{cos}}{\partial s} (s_j)
 =&\, \begin{cases}
        \frac{1}{\Delta s} \left( \hat{B}_{\theta, mn}^\mathrm{cos}(s_{j+0.5}) - \hat{B}_{\theta, mn}^\mathrm{cos}(s_{j-0.5}) \right) &:\, \textrm{ even } m \\
        \frac{\sqrt{s_j}}{\Delta s} \left(
            \hat{b}_{\theta, mn}^\mathrm{cos}(s_{j+0.5})
          - \hat{b}_{\theta, mn}^\mathrm{cos}(s_{j-0.5}) \right) \\
        ~ + \frac{1}{4 \sqrt{s_j}} \left(
              \hat{b}_{\theta, mn}^\mathrm{cos}(s_{j+0.5})
            + \hat{b}_{\theta, mn}^\mathrm{cos}(s_{j-0.5}) \right)
          &:\, \textrm{ odd } m
      \end{cases} \\
 \frac{\partial \hat{B}_{\varphi, mn}^\mathrm{cos}}{\partial s} (s_j)
 =&\, \begin{cases}
        \frac{1}{\Delta s} \left( \hat{B}_{\varphi, mn}^\mathrm{cos}(s_{j+0.5}) - \hat{B}_{\varphi, mn}^\mathrm{cos}(s_{j-0.5}) \right) &:\, \textrm{ even } m \\
        \frac{\sqrt{s_j}}{\Delta s} \left(
            \hat{b}_{\varphi, mn}^\mathrm{cos}(s_{j+0.5})
          - \hat{b}_{\varphi, mn}^\mathrm{cos}(s_{j-0.5}) \right) \\
        ~ + \frac{1}{4 \sqrt{s_j}} \left(
              \hat{b}_{\varphi, mn}^\mathrm{cos}(s_{j+0.5})
            + \hat{b}_{\varphi, mn}^\mathrm{cos}(s_{j-0.5}) \right)
          &:\, \textrm{ odd } m
      \end{cases}
\end{align}
with
\begin{align}
 \hat{b}_{\theta, mn}^\mathrm{cos}(s_{j+0.5}) =&\, \frac{\hat{B}_{\theta, mn}^\mathrm{cos}(s_{j+0.5})}{\sqrt{s_{j+0.5}}} \\
 \hat{b}_{\varphi, mn}^\mathrm{cos}(s_{j+0.5}) =&\, \frac{\hat{B}_{\varphi, mn}^\mathrm{cos}(s_{j+0.5})}{\sqrt{s_{j+0.5}}} \, .
\end{align}
%
This is the differencing scheme introduced in Ref.~\cite{hirshman_schwenn_nuehrenberg_1990}.
%
The regularized interpolation of~$\hat{B}_{s,mn}^\mathrm{sin}$ is done as follows:
\begin{align}
 \hat{B}_{s,mn}^\mathrm{sin}
 = \begin{cases}
    \frac{1}{2} \left( \hat{B}_{s, mn}^\mathrm{sin}(s_{j+0.5}) + \hat{B}_{s, mn}^\mathrm{sin}(s_{j-0.5}) \right) &:\, \textrm{ even } m \\
    \frac{1}{2 \sqrt{s_j}} \left(
        \hat{b}_{s, mn}^\mathrm{cos}(s_{j+0.5})
      + \hat{b}_{s, mn}^\mathrm{cos}(s_{j-0.5}) \right) &:\, \textrm{ odd } m
   \end{cases}
\end{align}
with
\begin{align}
 \hat{b}_{s, mn}^\mathrm{cos}(s_{j+0.5})
 = \sqrt{s_{j+0.5}} \hat{B}_{s,mn}^\mathrm{sin}(s_{j+0.5}) \, .
\end{align}
%
The mapping between intermediate quantities occuring throughout the computation of the current density in VMEC
and the corresponding variables in the Fortran VMEC implementation is found in Tab.~\ref{tab:compute_currents_vars}.
\begin{table}[tb]
 \centering
 \begin{tabular}{|c|c|}
  \hline
  Fortran Variable name & Physical quantity \\
  \hline
  \texttt{bu0} & $\hat{b}_{\theta, mn}^\mathrm{cos}(s_{j-0.5})$ \\
  \texttt{bu1} & $\hat{b}_{\theta, mn}^\mathrm{cos}(s_{j+0.5})$ \\
  \texttt{bv0} & $\hat{b}_{\varphi, mn}^\mathrm{cos}(s_{j-0.5})$ \\
  \texttt{bv1} & $\hat{b}_{\varphi, mn}^\mathrm{cos}(s_{j+0.5})$ \\
  \hline
  \texttt{t1} & $\hat{B}_{s,mn}^\mathrm{sin} (s_j)$ \\
  \texttt{t2} & $\frac{\partial \hat{B}_{\theta, mn}^\mathrm{cos}}{\partial s} (s_j)$ \\
  \texttt{t3} & $\frac{\partial \hat{B}_{\varphi, mn}^\mathrm{cos}}{\partial s} (s_j)$ \\
  \hline
 \end{tabular}
 \caption{Fortran variable names for the various physical quantities
          occuring throughout the computation of the current density in VMEC.}
 \label{tab:compute_currents_vars}
\end{table}
%
The \texttt{b\{u,v\}\{0,1\}} terms are used to compute the \texttt{t\{2,3\}} terms as follows:
\begin{align}
 \texttt{t2} =&\, \begin{cases}
                   \frac{   1      }{\Delta s} \left( \hat{B}_{\theta, mn}^\mathrm{cos}(s_{j+0.5}) - \hat{B}_{\theta, mn}^\mathrm{cos}(s_{j-0.5}) \right)     &:\, \textrm{even } m \\
                   \frac{\sqrt{s_j}}{\Delta s} \left( \texttt{bu1} - \texttt{bu0} \right) + \frac{1}{4 \sqrt{s_j}} \left( \texttt{bu1} + \texttt{bu0} \right) &:\, \textrm{odd  } m
                  \end{cases} \\
 \texttt{t3} =&\, \begin{cases}
                   \frac{   1      }{\Delta s} \left( \hat{B}_{\varphi, mn}^\mathrm{cos}(s_{j+0.5}) - \hat{B}_{\varphi, mn}^\mathrm{cos}(s_{j-0.5}) \right)       &:\, \textrm{even } m \\
                   \frac{\sqrt{s_j}}{\Delta s} \left( \texttt{bv1} - \texttt{bv0} \right) + \frac{1}{4 \sqrt{s_j}} \left( \texttt{bv1} + \texttt{bv0} \right) &:\, \textrm{odd  } m  \, .
                  \end{cases}
\end{align}
%
Using these definitions, the Fourier coefficients of the current density (times the Jacobian)
are computed as follows in the VMEC subroutine \texttt{Compute\_Currents}:
\begin{align}
 \texttt{currumnc} = (\sqrt{g} j^\theta)^{\textrm{cos}}_{mn}
   =&\, \frac{1}{\mu_0} \left( -n \cdot \texttt{t1} - \texttt{t3} \right) \\
 \texttt{currvmnc} = (\sqrt{g} j^\varphi )^{\textrm{cos}}_{mn}
   =&\, \frac{1}{\mu_0} \left( -m \cdot \texttt{t1} + \texttt{t2} \right) \, .
\end{align}
%
Linear extrapolation is used to compute the current density Fourier coefficents with $m \leq 1$ at the magnetic axis,
with all other Fourier harmonics set to zero at the magnetic axis.
%
Linear extrapolation is also used to compute the current density coefficients at the LCFS for all Fourier harmonics.
%
The source term in the Biot-Savart integral over the plasma current density,
namely the cylindrical components of the current density~$(j^{R}, j^{\varphi}, j^{Z})$ times the Jacobian,
are then computed from the contravariant representation as follows
for the remaining $j^R$ and $j^Z$:
\begin{align}
 \sqrt{g} j^{R}       =&\,   \sqrt{g} j^\theta \frac{\partial R}{\partial \theta} + \sqrt{g} j^\varphi \frac{\partial R}{\partial \varphi} \\
 \sqrt{g} j^{Z}       =&\,   \sqrt{g} j^\theta \frac{\partial Z}{\partial \theta} + \sqrt{g} j^\varphi \frac{\partial Z}{\partial \varphi} \, .
\end{align}

\section{Field Line Following in VMEC}
Straight-fieldline-coordinate are $(s, \theta^*, \zeta)$ in VMEC, where
\begin{equation}
 \theta^* = \theta + \lambda(s,\theta,\zeta) \label{eqn:theta_star}
\end{equation}
and $(s,\theta,\zeta)$ are the VMEC-internal coordinates
that lead to spectrally-condensed Fourier representations for $R$ and $Z$.
The stream function $\lambda$ is expressed as a Fourier series as well
(considering only the stellarator-symmetric case for now):
\begin{equation}
 \lambda(s,\theta,\zeta) = \sum\limits_{m,n} \hat{\lambda}^\textrm{sin}_{m,n}(s) \sin(m \theta - n \zeta) \, .
\end{equation}
Assume a certain location $(s,\theta,\zeta)$ is to be mapped along the magnetic field lines
to another toroidal location $\zeta'$.
In order to achieve this, we need to:
\begin{enumerate}
 \item compute $\theta^*_\textrm{target} = \theta^*(s,\theta,\zeta)$ and
 \item find $\theta'$ such that $\theta^*(s,\theta',\zeta') = \theta^*_\textrm{target}$.
\end{enumerate}
The first point is done by simply evaluating \eqn{theta_star}.
For the second step, a non-linear root-finding method can be employed to
find the root of the following objective function:
\begin{equation}
 f(\theta') = \theta^*(s,\theta',\zeta') - \theta^*_\textrm{target} \, . \label{eqn:theta_star_objective}
\end{equation}
It can be safely assumed that there exists only one solution
and furthermore, $\theta^*$ can be assumed to be monotonic with respect to $\theta$.
This allows to use a Newton method for solving \eqn{theta_star_objective}.
The Newton method needs derivate information, which can be computed analytically:
\begin{equation}
 \frac{\partial f}{\partial \theta}(s,\theta',\zeta')
 =
 \frac{\partial \theta^*}{\partial \theta}(s,\theta',\zeta')
 =
 1 + \frac{\partial \lambda}{\partial \theta}(s,\theta',\zeta') \, . \label{eqn:df_dthetap}
\end{equation}
The derivative of $\lambda$ can be computed analytically as well:
\begin{equation}
 \frac{\partial \lambda}{\partial \theta}(s,\theta',\zeta')
 =
 \sum\limits_{m,n} m \hat{\lambda}^\textrm{sin}_{m,n}(s) \cos(m \theta' - n \zeta') \, .
\end{equation}
It is expected that the objective function $f$ needs to be evaluated multiple times
for different values of $\theta'$, but always at the given values of $s$ and $\zeta'$
during the root-finding procedure.
Thus, a transformed Fourier representation for $\lambda$ and $\partial \lambda/\partial \theta$
can be derived:
\begin{align}
 \lambda\vert_{\zeta'}(\theta')
 =&\,
 \sum\limits_{m,n} \hat{\lambda}^\textrm{sin}_{m,n}(s) \sin(m \theta' - n \zeta') \nonumber \\ ~
 =&\,
 \sum\limits_{m} \sum\limits_{n} \hat{\lambda}^\textrm{sin}_{m,n}(s) \left[
   \sin(m \theta') \cos(n \zeta') - \cos(m \theta') \sin(n \zeta')
 \right] \nonumber \\ ~
 =&\,
 \sum\limits_{m} \left[
   \Bigl(\sum\limits_{n} \hat{\lambda}^\textrm{sin}_{m,n}(s) \cos(n \zeta') \Bigr) \sin(m \theta') -
   \Bigl(\sum\limits_{n} \hat{\lambda}^\textrm{sin}_{m,n}(s) \sin(n \zeta') \Bigr) \cos(m \theta')
 \right] \nonumber \\ ~
 =&\,
 \sum\limits_{m} \left[
   \hat{\lambda}^\textrm{sin}_{m,\zeta'}(s) \sin(m \theta') -
   \hat{\lambda}^\textrm{cos}_{m,\zeta'}(s) \cos(m \theta')
 \right] \label{eqn:lambda_at_zetap}
\end{align}
with
\begin{align}
 \hat{\lambda}^\textrm{sin}_{m,\zeta'}(s)
 =&\, \sum\limits_{n} \hat{\lambda}^\textrm{sin}_{m,n}(s) \cos(n \zeta') \label{eqn:lmszeta} \\
 \hat{\lambda}^\textrm{cos}_{m,\zeta'}(s)
 =&\, \sum\limits_{n} \hat{\lambda}^\textrm{sin}_{m,n}(s) \sin(n \zeta') \label{eqn:lmczeta} \, .
\end{align}
The coefficients $\hat{\lambda}^\textrm{sin}_{m,\zeta'}(s)$ and $\hat{\lambda}^\textrm{cos}_{m,\zeta'}(s)$
can be computed once for the given values of $s$ and $\zeta'$
and then only the summation over $m$ in \eqn{lambda_at_zetap} has to be performed
in every iteration of the root-finding algorithm.
%
Similarly, the computation of $\partial \lambda/\partial \theta$ can be streamlined,
now already based on \eqn{lambda_at_zetap}:
\begin{equation}
 \frac{\partial \lambda}{\partial \theta}\vert_{\zeta'}(\theta')
 =
 \sum\limits_{m} \left[
   m \hat{\lambda}^\textrm{sin}_{m,\zeta'}(s) \cos(m \theta') +
   m \hat{\lambda}^\textrm{cos}_{m,\zeta'}(s) \sin(m \theta')
 \right] \label{eqn:lu_at_zetap}
\end{equation}
Thus, the following algorithm can be used for field-line following on VMEC equilibria:
\begin{enumerate}
 \item Assume that $s$, $\theta$, $\zeta$ and $\zeta'$ are given.
       Assume that the Fourier coefficients $\hat{\lambda}^\textrm{sin}_{m,n}(s)$ are given for all $m$ and $n$.
       The desired tolerance is $\epsilon$.
 \item Compute $\theta^*_\textrm{target}$ using \eqn{theta_star}.
 \item For each $m$, pre-compute \eqn{lmszeta} and \eqn{lmczeta}.
 \item Initialize $\theta' \leftarrow \theta$ as an initial guess.
 \item Evaluate the objective function \eqn{theta_star_objective}:
       \begin{enumerate}
         \item Compute $\lambda\vert_{\zeta'}(\theta')$ using \eqn{lambda_at_zetap}.
         \item Compute $\theta^*(\theta')$ using \eqn{theta_star}.
         \item Compute $f(\theta')$ using \eqn{theta_star_objective}.
       \end{enumerate}
 \item Consider converged if $f(\theta') < \epsilon$ and exit if converged.
 \item Compute the derivative of the objective function \eqn{df_dthetap}:
       \begin{enumerate}
         \item Compute $\partial \lambda / \partial \theta \vert_{\zeta'}(\theta')$ using \eqn{lu_at_zetap}.
         \item Compute $\partial f / \partial \theta\vert_{\zeta'}(\theta')$ using \eqn{df_dthetap}.
       \end{enumerate}
 \item Perform the Newton step:
       \begin{equation}
         \theta' \leftarrow \theta' - \frac{f(\theta')}{\partial f / \partial \theta (\theta')} \, . \nonumber
       \end{equation}
 \item Go back to step 5.
\end{enumerate}
The result of this algorithm is a set of VMEC coordinates $(s,\theta',\zeta')$,
which are on the same field line as the VMEC coordinates $(s,\theta,\zeta)$.

\section{Inverse Coordinate Transform}
%
Assume flux surface shapes are given in VMEC(-like) coordinates $R(s, \theta, \varphi)$ and $Z(s, \theta, \varphi)$
as function of normalized toroidal flux $s$ with radius-like radial coordinate $\rho = \sqrt{s}$
and angle-like poloidal coordinate $\theta$ and toroidal coordinate $\varphi$,
equal to the cylindrical coordinates angle.
%
The representation is assumed to be as Fourier series for a number of discrete flux surfaces:
\begin{align}
 R(s_j, \theta, \varphi) =&\, \sum_{m,n} R^\textrm{cos}_{mn}(s_j) \cos(m \theta - n n_\textrm{fp} \varphi) \label{eqn:r_of_flux_coords} \\
 Z(s_j, \theta, \varphi) =&\, \sum_{m,n} Z^\textrm{sin}_{mn}(s_j) \sin(m \theta - n n_\textrm{fp} \varphi) \label{eqn:z_of_flux_coords}
\end{align}
with the number of toroidal field periods $n_\textrm{fp}$
at the $j$-th flux surface, identified by radial coordinate $s_j$.
%
A point in space (inside or outside the plasma) is identified by cylindrical coordinates $(R_0, \varphi_0, Z_0)$.
%
The goal of the method described here is to find VMEC coordinates $(s, \theta, \varphi)$,
for a given equilibrium and for now, assuming that the point under consideration is within the plasma volume,
corresponding to the point's real-space coordinates $(R_0, \varphi_0, Z_0)$.

\FloatBarrier
\subsection{Basic Approach}
%
It is noted that there is no direct inverse of the coordinate parameterization
in \eqn{r_of_flux_coords} and \eqn{z_of_flux_coords}.
%
In turn, the approach used here is to formulate this inverse coordinate transform as a non-linear minimization problem:
\begin{align}
 \min_{0 \leq s \leq 1, 0 \leq \theta < 2 \pi} \delta(s, \theta)
\end{align}
with the real-space distance~$\delta$ between the trial point $(R(s, \theta, \varphi_0), \varphi_0, Z(s, \theta, \varphi_0))$
and the given point $(R_0, \varphi_0, Z_0)$, whose VMEC coordinates~$(s, \theta)$ are to be inferred:
\begin{equation}
 \delta^2 (s, \theta) = \left[ R(s, \theta, \varphi_0) - R_0 \right]^2 + \left[ Z(s, \theta, \varphi_0) - Z_0 \right]^2 \, .
\end{equation}
%
This implies that the minimization has to be carried out only in the poloidal cutplane over $(s, \theta)$
at a fixed toroidal angle $\varphi_0$ for the given point.

\FloatBarrier
\subsection{Initial Guess}
%
A good initial guess $(s^{(-1)}, \theta^{(-1)})$ can greatly accelerate convergence of this iterative method.
%
Assume the magnetic axis location at toroidal angle $\varphi_0$ is given by $(R_\textrm{axis}(\varphi_0), Z_\textrm{axis}(\varphi_0))$ (see below).
Then, the distance to the target point is given by $d$:
\begin{align}
 \Delta R_\textrm{axis} =&\, R_0 - R_\textrm{axis}(\varphi_0) \\
 \Delta Z_\textrm{axis} =&\, Z_0 - Z_\textrm{axis}(\varphi_0) \\
 d =&\, \sqrt{ \Delta R_\textrm{axis}^2 + \Delta Z_\textrm{axis}^2 } \, .
\end{align}
%
Assume the minor radius of the plasma is approximated by $a_\textrm{eff}$ (see below).
%
Then, the initial guess for the radial coordinate can be approximated as $\rho^{(-1)} = d / a_\textrm{eff}$.
%
The initial guess for the poloidal coordinate~$\theta^{(-1)}$ is computed
by computing the simple poloidal angle of the point in question
with respect to the magnetic axis:
\begin{align}
 \theta^{(-1)} = \texttt{atan2}(\Delta Z_\textrm{axis}, \Delta R_\textrm{axis}) \, .
\end{align}

\subsubsection{Magnetic Axis Geometry}
The magnetic axis is the central flux surface, which collapses to a filament.
%
Its geometry is given as a purely toroidal Fourier series,
since the magnetic axis filament has no poloidal dependence:
\begin{align}
 R_\textrm{axis}(\varphi) = \sum_{n=0}^{\texttt{ntor}} R^\textrm{cos}_{0n}(0) \cos(n n_\textrm{fp} \varphi) \label{eqn:axis_r} \\
 Z_\textrm{axis}(\varphi) = \sum_{n=0}^{\texttt{ntor}} Z^\textrm{sin}_{0n}(0) \sin(n n_\textrm{fp} \varphi) \label{eqn:axis_z} \, .
\end{align}

\subsubsection{Minor Radius Approximation}
%
In the simplest case, where only the $(m=1,n=0)$ harmonic is non‐zero,
the cross‐section is an ellipse in the poloidal plane:
\begin{align}
 R(\theta) =&\, R^\textrm{cos}_{0,0} + R^\textrm{cos}_{1,0} \cos(\theta) \\
 Z(\theta) =&\,                        Z^\textrm{sin}_{1,0} \sin(\theta)
\end{align}
%
If $R^\textrm{cos}_{1,0} = Z^\textrm{sin}_{1,0} = a$, that is a perfect circle of radius $a$.
%
If $R^\textrm{cos}_{1,0} \neq Z^\textrm{sin}_{1,0}$, you get an ellipse with semi-axes $R^\textrm{cos}_{1,0}$ and $Z^\textrm{sin}_{1,0}$.
%
The geometric mean of those semi-axes
\begin{equation}
 a_\textrm{eff} = \sqrt{R^\textrm{cos}_{1,0} \cdot Z^\textrm{sin}_{1,0}} \label{eqn:minor_radius_approximation}
\end{equation}
is exactly the radius of a circle having the same area~$A$ as that ellipse
($A = \pi R^\textrm{cos}_{1,0} Z^\textrm{sin}_{1,0}$, so $\pi a_\textrm{eff}^2 = A$).
In plasma physics, the minor radius is often taken to be the radius of a circle of equal cross‐sectional area
as the flux surface under consideration.

\FloatBarrier
\subsection{Iterative Algorithm}
%
A Newton-type iteration algorithm is employed in this context~\cite{attenberger_1987}.
%
For brevity, we consider in the following only $R(s,\theta)$ and $Z(s,\theta)$.
%
A Taylor expansion can be used to approximate $R_0$ and $Z_0$ around $s^{(k)}$ and $\theta^{(k)}$
in the $k$-th iteration:
\begin{align}
 R_0 =&\, R(s^{(k)},\theta^{(k)}) +
          (     s -      s^{(k)}) \frac{\partial R}{\partial      s}(s^{(k)}, \theta^{(k)}) +
          (\theta - \theta^{(k)}) \frac{\partial R}{\partial \theta}(s^{(k)}, \theta^{(k)}) + ... \label{eqn:approx_r} \\
 Z_0 =&\, Z(s^{(k)},\theta^{(k)}) +
          (     s -      s^{(k)}) \frac{\partial Z}{\partial      s}(s^{(k)}, \theta^{(k)}) +
          (\theta - \theta^{(k)}) \frac{\partial Z}{\partial \theta}(s^{(k)}, \theta^{(k)}) + ... \label{eqn:approx_z} \, .
\end{align}
The following abbreviations are introduced:
\begin{align}
 R^{(k)}        \equiv&\, R(s^{(k)},\theta^{(k)}) \\
 R^{(k)}_s      \equiv&\, \frac{\partial R}{\partial      s}(s^{(k)}, \theta^{(k)}) \\
 R^{(k)}_\theta \equiv&\, \frac{\partial R}{\partial \theta}(s^{(k)}, \theta^{(k)})
\end{align}
and equivalently for $Z$.
%
If the Taylor expansions in \eqn{approx_r} and \eqn{approx_z} are terminated after the shown terms,
these expressions form a rank-2 linear system of equations:
\begin{equation}
  \begin{pmatrix}
    R_0 - R^{(k)} \\
    Z_0 - Z^{(k)}
  \end{pmatrix}
  =
  \begin{pmatrix}
    R^{(k)}_s & R^{(k)}_\theta \\
    Z^{(k)}_s & Z^{(k)}_\theta
  \end{pmatrix}
  \begin{pmatrix}
    s - s^{(k)} \\
    \theta - \theta^{(k)}
  \end{pmatrix} \label{eqn:approx_rz} \, .
\end{equation}
This system can be solved for $(s - s^{(k)})$ and $(\theta - \theta^{(k)})$.
The determinant of the matrix in \eqn{approx_rz} is:
\begin{equation}
 \tau^{(k)} = R^{(k)}_s Z^{(k)}_\theta - R^{(k)}_\theta Z^{(k)}_s \, .
\end{equation}
Then, the solution is given by:
\begin{equation}
 \begin{pmatrix}
         s \\
    \theta
  \end{pmatrix}
  =
 \begin{pmatrix}
         s^{(k)} \\
    \theta^{(k)}
  \end{pmatrix}
  +
  \frac{1}{\tau^{(k)}}
  \begin{pmatrix}
     Z^{(k)}_\theta & -R^{(k)}_\theta \\
    -Z^{(k)}_s      &  R^{(k)}_s
  \end{pmatrix}
  \begin{pmatrix}
    R_0 - R^{(k)} \\
    Z_0 - Z^{(k)}
  \end{pmatrix} \, .
\end{equation}
This can be used as an iteration equation for $s$ and $\theta$,
where the values for the $(k+1)$-th iteration are given by:
\begin{align}
 s^{(k+1)} =&\, s^{(k)} +
   \frac{ Z^{(k)}_\theta \left( R_0 - R^{(k)} \right) -
          R^{(k)}_\theta \left( Z_0 - Z^{(k)} \right) }{\tau^{(k)}} \\
 \theta^{(k+1)} =&\, \theta^{(k)} +
   \frac{ R^{(k)}_s \left( Z_0 - Z^{(k)} \right) -
          Z^{(k)}_s \left( R_0 - R^{(k)} \right) }{\tau^{(k)}} \, .
\end{align}
The inverse coordinate transform is considered converged
if the distance between the targeted point $(R_0, Z_0)$ and the model values $(R^{(k)},Z^{(k)})$
is smaller than a prescribed tolerance~$\epsilon$:
\begin{equation}
 \left( R_0 - R^{(k)} \right)^2 + \left( Z_0 - Z^{(k)} \right)^2 < \epsilon^2 \, .
\end{equation}
In case the iteration gives a worse result than the previous iterate,
the step size is halved
and an averaged value of the Jacobian is used:
\begin{align}
 s^{(k+2)} =&\, s^{(k)} + \frac{1}{2}
   \frac{ Z^{(k)}_\theta \left( R_0 - R^{(k)} \right) -
          R^{(k)}_\theta \left( Z_0 - Z^{(k)} \right) }{0.75 \tau^{(k)} + 0.25 \tau^{(k+1)}} \label{eqn:adaptive_step_s} \\
 \theta^{(k+2)} =&\, \theta^{(k)} + \frac{1}{2}
   \frac{ R^{(k)}_s \left( Z_0 - Z^{(k)} \right) -
          Z^{(k)}_s \left( R_0 - R^{(k)} \right) }{0.75 \tau^{(k)} + 0.25 \tau^{(k+1)}} \label{eqn:adaptive_step_u} \, .
\end{align}
This process can be repeated as necessary
until a new iterate for $s$ and $\theta$ is found
which reduces the objective.

\FloatBarrier
\subsection{Radial Derivative of VMEC Flux Surface Geometry}
%
One key point to consider when working with flux surface geometries from VMEC
is the interpolation of the flux surface geometry and quantities on them
in the radial direction.
%
This is because VMEC solves only for discrete flux surfaces
and uses finite differences in the radial direction,
which leaves it to the user to choose how to interpolate quantities between flux surfaces.
%
Here, an interpolation approach is described, which was found to work very well in general,
and in particular for the coordinate transform use case considered here.

The interpolation is set up over $\rho = \sqrt{s}$ as the radial coordinate.
%
Cubic splines are used with natural boundary conditions to set up the actual interpolation.
%
The interpolation and associated analytical derivatives are set up
for the radial dependence of each Fourier coefficient in the representations
\eqn{r_of_flux_coords} and \eqn{z_of_flux_coords} individually.
%
The boundary condition at the plasma boundary is the natural boundary condition provided by the spline implementation.
%
The boundary condition at the magnetic axis, i.e., at the center of the plasma,
is adjusted to satisfy the poloidal parity of the respective Fourier coefficients to interpolate:
\begin{align}
 X_{m, n}(-s) =&\, \phantom{-} ~ X_{m, n}(s) \textrm{ if } m \textrm{ even } \\
 X_{m, n}(-s) =&\,          -    X_{m, n}(s) \textrm{ if } m \textrm{ odd }
\end{align}
for $X \in \{ R^\textrm{cos}, Z^\textrm{sin} \}$.
%
Importantly, note that the parity of the interpolation at the magnetic axis
only depends on the poloidal mode number and not on the phase (sin or cos)
of the function to interpolate.
%
This regularization is particularly relevant for obtaining reliable
inverse coordinate transforms near the magnetic axis, i.e.,
inside of the first flux surface of the VMEC equilibrium under consideration.

The geometry of the flux surfaces in VMEC is represented on the ``full'' radial grid:
\begin{equation}
 s_j = \frac{j}{\texttt{ns} - 1} \textrm{ for } j = 0, 1, ..., (\texttt{ns} - 1)
\end{equation}
where $j=0$ represents the magnetic axis, i.e., the core of the plasma
and $j=\texttt{ns} - 1$ represents the plasma boundary.

The even and odd symmetry boundary conditions at the left end of the range
in the radial direction (at the magnetic axis) can be enforced in some spline
interpolation implementations.
%
Some spline interpolation implementations
feature only natural boundary conditions at both ends.
%
One can mirror the data according to the desired symmetry
and interpolate over that extended dataset
and thereby still establish symmetry with the desired parity at the magnetic axis
using such a spline implementation.
%
The following concrete formulation is used
for even symmetry at the magnetic axis when interpolating Fourier coefficients on the full-grid:
\begin{align}
 \rho[\texttt{ns} - 1 + j] =&\, \sqrt{ \frac{j}{\texttt{ns} - 1} } \textrm{ for } 0 \leq j < \texttt{ns} \\
 f[\texttt{ns} - 1 + j]    =&\, X[j]                               \textrm{ for } 0 \leq j < \texttt{ns} \\
 \textrm{ and }             & \textrm{ then } \nonumber \\
 \rho[\texttt{ns} - 1 - j] =&\, -\rho[\texttt{ns} - 1 + j]         \textrm{ for } 1 \leq j < \texttt{ns} \\
 f[\texttt{ns} - 1 - j]    =&\, f[\texttt{ns} - 1 + j]             \textrm{ for } 1 \leq j < \texttt{ns}
\end{align}
where $\rho$ and $f$ are the arrays of length $2 \cdot \texttt{ns} - 1$ to be handed to the spline interpolation
and $X$ is the Fourier coefficient array on the full grid to be interpolated.
%
For enforcing odd parity at the magnetic axis,
the following formulation is used instead:
\begin{align}
 \rho[\texttt{ns} - 1 + j] =&\, \sqrt{ \frac{j}{\texttt{ns} - 1} } \textrm{ for } 0 \leq j < \texttt{ns} \\
 f[\texttt{ns} - 1 + j]    =&\, X[j]                               \textrm{ for } 0 \leq j < \texttt{ns} \\
 \textrm{ and }             & \textrm{ then } \nonumber \\
 \rho[\texttt{ns} - 1 - j] =&\, -\rho[\texttt{ns} - 1 + j]         \textrm{ for } 1 \leq j < \texttt{ns} \\
 f[\texttt{ns} - 1 - j]    =&\, -f[\texttt{ns} - 1 + j]            \textrm{ for } 1 \leq j < \texttt{ns} \label{eqn:mirror_odd}
\end{align}
where the only difference to the formulation for even parity
is the negative sign in mirroring the data in \eqn{mirror_odd}.

\FloatBarrier
\subsection{Extrapolation beyond the Plasma Boundary}
%
VMEC only provides the flux surface geometry up to the plasma boundary.
%
Extrapolation of the geometry can be set up in a way that makes
the extrapolated flux surface geometry to tend towards ellipses far away from the plasma.
%
This works for most plasma boundary shapes, but note that it is possible
that the extrapolation scheme outlined below leads to extrapolated flux surfaces overlapping
with the VMEC-provided flux surfaces, in particular close to the plasma boundary.
%
The extrapolation method is described in Ref.~\cite{attenberger_1987}.
%
The Fourier coefficients of the flux surface geometry are extrapolated
to $\rho > 1$ as follows:
\begin{align}
 X_{mn}(\rho) =&\, \begin{cases}
                    X_{mn}(1)          \cdot \rho  & \textrm{ for } (m,n)=(1,0) \\
                    X_{mn}(1) \phantom{\cdot \rho} & \textrm{ else }
                   \end{cases}
\end{align}
for $X \in \{ R^\textrm{cos}, Z^\textrm{sin} \}$.
%
Therefore, the radial derivatives are extrapolated as follows
outside of the plasma boundary at $\rho > 1$:
\begin{align}
 \frac{\partial X_{mn}}{\partial \rho}(\rho) =&\, \begin{cases}
                                                   X_{mn}(1) & \textrm{ for } (m,n)=(1,0) \\
                                                   0         & \textrm{ else. }
                                                  \end{cases}
\end{align}

\FloatBarrier
\subsection{Actual Implementation of Inverse Coordinate Transform}
%
The actual implementation of the inverse coordinate transform for VMEC coordinates
is suggested to always use the extrapolation for $\rho > 1$ introduced above.
%
If points outside the plasma boundary shall be marked as such,
one can identify cases where $\rho > 1$ was found as radial VMEC coordinate
for a given real-space point, and replace the inferred coordinates
with $\hat{\rho} = \infty$ and $\hat{\theta} = \textrm{NaN}$, $\hat{\varphi} = \textrm{NaN}$.
%
Note that the actual implementation uses $\rho$ as the radial coordinate,
since this was found to work better with the interpolation
of flux surface geometry near the magnetic axis.
%
The actual implementation is presented in Algorithm~\ref{alg:inverse_coordinate_transform}.
%
\begin{algorithm}
  \caption{Inverse Coordinate Transform $(R_0, \varphi_0, Z_0) \to (\rho, \theta, \varphi_0)$ with extrapolation}
  \begin{algorithmic}[1]
    \Require
      $R_0, \varphi_0, Z_0$ ; $\epsilon > 0$ ; $n_\textrm{max} \in \mathbb{N}$

    % Initialize squared tolerance
    \State $\epsilon^2 \gets \epsilon \cdot \epsilon$

    % Compute magnetic axis and distance to axis
    \State $\{ R_\textrm{axis}(\varphi_0), Z_\textrm{axis}(\varphi_0) \} \gets \texttt{GetMagneticAxis}(\varphi_0)$
    \State $\Delta R_\textrm{axis} \gets R_0 - R_\textrm{axis}(\varphi_0)$ ; $\Delta Z_\textrm{axis} \gets Z_0 - Z_\textrm{axis}(\varphi_0)$
    \State $d \gets \sqrt{\Delta R_\textrm{axis}^2 + \Delta Z_\textrm{axis}^2}$

    % Approximate minor radius and initial guess
    \State $r_0 \gets \texttt{MinorRadiusApproximation}()$
    \State $ \rho \gets d / r_0$
    \State $\theta \gets \texttt{atan2}(\Delta Z_\textrm{axis}, \Delta R_\textrm{axis})$
    \If{$\theta < 0$}
      \State $\theta \gets \theta + 2\pi$
    \EndIf

    % Initialize iteration variables
    \State $\delta^2 \gets +\infty$ ; ${\delta^2}^{(-1)} \gets +\infty$
    \State $\beta \gets 1$
    \State $\tau^{(-1)} \gets \textrm{NaN}$

    \For{$i = 1$ to $n_\textrm{max}$}

      % Evaluate forward map and residuals
      \State Evaluate $R$, $\partial R / \partial \rho$ and $\partial R / \partial \theta$ at $(\rho, \theta)$ and analogously for $Z$
      \State $\Delta R \gets R(\rho, \theta) - R_0$ ; $\Delta Z \gets Z(\rho, \theta) - Z_0$
      \State ${\delta^2}^{(-1)} \gets \delta^2$ ; $\delta^2 \gets \Delta R^2 + \Delta Z^2$
      \If{$\delta^2 < \epsilon^2$}
        \State \textbf{break}
      \EndIf

      % Compute Jacobian determinant
      \State $\tau^{(-1)} \gets \tau$
      \State $\tau \gets \partial R / \partial \rho (\rho, \theta) \cdot \partial Z / \partial \theta (\rho, \theta) - \partial R / \partial \theta (\rho, \theta) \cdot \partial Z / \partial \rho (\rho, \theta)$
      \If{$i=1$}
        \State $\tau^{(-1)} \gets \tau$
      \EndIf

      % Handle case of un-successful step - distance to target increased -> mix Jacobian with previous Jacobian and reduce step length
      \If{$\delta^2 > {\delta^2}^{(-1)}$}
        \State $\tau \gets 0.75\,\tau^{(-1)} + 0.25\,\tau$ ; $\beta \gets \beta/2$
      \Else
        % Restore original step length in case of a successful step
        \State $\beta \gets 1$
      \EndIf

      % Compute next iterate via Newton-like update
      \State $\rho   \gets \rho   - \beta \left[\partial Z / \partial \theta (\rho, \theta) \cdot \Delta R - \partial R / \partial \theta (\rho, \theta) \cdot \Delta Z \right] / \tau$
      \State $\theta \gets \theta - \beta \left[\partial R / \partial \rho   (\rho, \theta) \cdot \Delta Z - \partial Z / \partial \rho   (\rho, \theta) \cdot \Delta R \right] / \tau$

      % Handle crossing over axis
      \If{$\rho < 0$}
        \State $\theta \gets \theta + \pi$ ; $\rho \gets -\rho$
      \EndIf

      % Make sure \theta stays within [0, 2 pi]
      \While{$\theta < 0$}
        \State $\theta \gets \theta + 2 \pi$
      \EndWhile
      \State $\theta \gets \theta \bmod 2 \pi$

    \EndFor

    \State \Return $(\rho, \theta, \varphi_0)$
  \end{algorithmic}
  \label{alg:inverse_coordinate_transform}
\end{algorithm}
%
Note that the step is taken in the ``reverse'' direction,
which is probably required to incorporate the negative Jacobian used in the VMEC coordinate system.
%
The method $\texttt{GetMagneticAxis}(\varphi_0)$ implements \eqn{axis_r} and \eqn{axis_z}.
%
The method $\texttt{MinorRadiusApproximation}()$ implements \eqn{minor_radius_approximation}.

\FloatBarrier
\subsection{Testing Strategy}
%
The functionality provided by this inverse coordinate algorithm is actually quite fundamental
to stellarator modeling and data analysis.
%
It is therefore used in quite a number of places of the overall stellarator toolchain.
%
Therefore, this algorithm deserves significant, robust testing
in order to get an implementation ready for production use.

%
This section outlines a testing strategy that has proven sufficiently thorough
to get an implementation ready for large-scale use in data analysis for Wendelstein~7-X.
%
Credit for thinking out this testing strategy goes to Oliver Ford at IPP Greifswald.
%
This algorithm should be tested on a number of different VMEC equilibria,
which capture typical plasma shapes occuring in the context
where this coordinate transform algorithm is supposed to be used.
%
Random number generation is used to come up with test point coordinates.
%
It is important for debuggability to ensure proper random seed management
in the testing harness in order to be able to debug errors which occur also very late in a testing cycle,
which would be otherwise very hard to reproduce.
%
The tests are carried out it batches of points,
which each are generated with their own random seed,
which is drawn from a random number generator themselves (seeded with the ``seed of seeds'').
%
This multi-layer random number generation allows to reproduce errors
in a given batch of test points quicker than if all points are tested in a single sequence,
since the random number generator does not have to go through generating
all point coordinates starting from a given seed,
but only needs to go through number in given batch
to reproduce an error that occured in that batch.
%
Unit testing of the implementation is done in 10 batches of 1000 points each.
%
Extended testing to catch rare errors and really ensure that an implementation
is ready for production use should run at least 10000 batches of 5000 test points each.

Radial derivatives of the flux surface geometry are inaccurate near the plasma boundary,
due to the (assumed chosen) natural boundary condition for the spline there.
%
Therefore, the radial displacement of test points near the plasma boundary
is done via the tangent vector in poloidal direction,
whose components $(\partial R / \partial \theta, \partial Z / \partial \theta)$
can be computed analytically from the Fourier series representation of $R$ and $Z$.
%
This tangent vector is then rotated by 90 degrees into the perpendicular orientation
by swapping its components and inverting one of them (see below).

The inverse coordinate transform algorithm presented here
is an iterative method with a user-prescribed (but usually chosen statically for a given application) convergence tolerance.
%
This means that it cannot be expected that the real space point
corresponding to inferred VMEC coordinates is closer to the target point in distance
than the prescribed convergence tolerance.
%
It is therefore imperative to avoid false negatives during testing of this algorithm
to compute the real-space distance between the point corresponding to the inferred coordinates
and the target point.
%
More importantly, if a test is made to check if a point is tested that was supposed to be inside the plasma boundary
actually is inferred to be inside the plasma boundary,
the test point has to be placed at least as far away from the plasma boundary towards the inside of the plasma
as the convergence tolerance distance is specified in the coordinate transform.
%
Otherwise, the inferred point may be within the convergence tolerance radius
around the target point, but fall just outside the plasma boundary,
even though the original point was inside the plasma boundary,
and an unsuitably-setup test would fail.

The inverse coordinate transform algorithm is tested then on a few different cases:
\begin{enumerate}
 \item points equally distributed in $(\rho, \theta, \varphi)$ space, i.e., anywhere within the plasma volume;
 \item points close to the magnetic axis
       - this particularly stresses the coordinate interpolation near the magnetic axis;
 \item points close the plasma boundary, but still predictably within the plasma boundary; and
 \item points close the plasma boundary, but already predictably outside of the plasma boundary
       - this particularly tests the coordinate extrapolation.
\end{enumerate}
%
For each of these cases, real-space cylindrical coordinates of test points~$[R_0, \varphi_0, Z_0]$
are made available. For points sampled close to the plasma boundary (cases 3 and 4),
an additional flag~\texttt{outside\_flag} is returned, which indicates if a given point is supposed
to be located inside the plasma boundary or not.

The testing procedure then is orchestrated as follows:
\begin{enumerate}
 \item Sample random test point coordinates~$[R_0, \varphi_0, Z_0]$ according to a method randomly selected from above list of test cases.
 \item Use the inverse coordinate transform algorithm under test to obtain the corresponding VMEC coordinates~$[\rho, \theta, \varphi]$.
 \item Transform the obtained VMEC coordinates back into cylindrical coordinates~$[R, \varphi, Z]$.
 \item Test if the returned value for $\varphi$ matches $\varphi_0$ - they should be bit-identical.
 \item If the extrapolation is inhibited by replacing the VMEC coordinates of points outside of the plasma boundary
       with $\rho = +\infty$, it is checked that a provided \texttt{outside\_flag} (see below) is consistent with the point
       actually being inferred as inside or outside the plasma boundary.
 \item Compute the Eucledian distance between $(R_0, Z_0)$ and $(R, Z)$ and make sure it is below the convergence tolerance $\epsilon$.
       If not, consider the test failed.
 \item Repeat all steps above for every test case in the current batch.
 \item Repeat all steps above for every batch in the whole set of test points.
\end{enumerate}
%
The random generation of test points for each of the four cases is presented next:
%
\subsubsection{Points Equally Distributed in the Plasma Volume}
%
The points equally distributed in the plasma volume are generated as follows:
\begin{align}
  %
  % sample normalized flux coordinate in [0,1]
  s &\sim \mathcal{U}(0,1)\\
  %
  % sample poloidal angle in [0,2\pi]
  \theta &\sim \mathcal{U}(0,2\pi)\\
  %
  % sample toroidal angle in [0,2\pi]
  \varphi_0 &\sim \mathcal{U}(0,2\pi)\\
  %
  % map to cylindrical coordinates via the inverse transform
  (R_0, \varphi_0, Z_0) &= \texttt{ToCylinderCoordinates}\bigl(s, \theta, \varphi_0 \bigr)
\end{align}
%
where \texttt{ToCylinderCoordinates} implements \eqn{r_of_flux_coords} and \eqn{z_of_flux_coords}.


\subsubsection{Points Close to the Magnetic Axis}
%
Test points close to the magnetic axis are generated as follows:
\begin{align}
  %
  % approximate minor radius
  a_\textrm{eff} &= \texttt{MinorRadiusApproximation}() \\
  %
  % small reference distance along minor radius
  d &= \frac{a_\textrm{eff}}{10} \\
  %
  % random radial and vertical offsets in [0,d]
  dr &\sim \mathcal{U}(0,d) \\
  dz &\sim \mathcal{U}(0,d) \\
  %
  % random toroidal angle in [0,2π]
  \varphi_0 &\sim \mathcal{U}(0,2\pi) \\
  %
  % magnetic axis coordinates at this toroidal angle
  \bigl[R_{\mathrm{axis}}(\varphi_0),\,Z_{\mathrm{axis}}(\varphi_0)\bigr] &= \texttt{GetMagneticAxis}(\varphi_0) \\
  %
  % construct test point near the axis
  R_{0} &= R_{\mathrm{axis}}(\varphi) + dr \\
  Z_{0} &= Z_{\mathrm{axis}}(\varphi) + dz \, .
\end{align}


\subsubsection{Points Close to the Plasma Boundary}

\begin{align}
%
% random poloidal angle on boundary
\theta_0 &\sim \mathcal{U}(0,2\pi)\\
%
% random toroidal angle on boundary
\varphi_0 &\sim \mathcal{U}(0,2\pi)\\
%
% signed offset distance from boundary (≥ requested accuracy)
d & = \epsilon \cdot \bigl(\mathcal{U}(0,1) + 1\bigr)\\
%
% cylindrical coords on boundary surface s=1
\bigl[R_{b},\,\varphi_{b},\,Z_{b}\bigr] &= \texttt{ToCylinderCoordinates}\bigl(1, \theta_0, \varphi_0 \bigr)\\
%
% arc‐length element along boundary curve
\ell & = \sqrt{\bigl(dR/d\theta\bigr)^{2} + \bigl(dZ/d\theta\bigr)^{2}}\\
%
% perpendicular displacements of magnitude d
dR_{\perp} & = \frac{dZ/d\theta}{\ell}\,d\\
%
dZ_{\perp} & = -\,\frac{dR/d\theta}{\ell}\,d\\
%
% final test point and outside flag (50\% inside, 50\% outside)
\bigl[R_{0},\,\varphi_{0},\,Z_{0},\,\texttt{outside\_flag}\bigr] &=
\begin{cases}
\bigl[R_{b}-dR_{\perp},\,\varphi_{b},\,Z_{b}-dZ_{\perp},\, \texttt{false} \bigr] \textrm{ in 50 \% of cases }  \\
\bigl[R_{b}+dR_{\perp},\,\varphi_{b},\,Z_{b}+dZ_{\perp},\, \texttt{true } \bigr] \textrm{ in 50 \% of cases }
\end{cases}
\end{align}
%
where $(dR_{\perp}, dZ_{\perp})$ is the outward-facing unit vector on the plasma boundary
and \texttt{outside\_flag} indicates to the test harness if a given test point
is supposed to be inside the plasma boundary or outside of it.

% \section{Computing the Magnetic Field}
% This section introduces how to compute the magnetic field from a VMEC equilibrium.
% Several methods are described here:
% \begin{enumerate}
%  \item Computation of the contravariant magnetic field components
%        from $\lambda$, $\iota$, $\sqrt{g}$ and $\Phi'$
%  \item Computation of the cylindrical magnetic field components
%        from the contravariant magnetic field components
%  \item Computation of the Cartesian magnetic field components
%        from the cylindrical magnetic field components.
% \end{enumerate}


% {\color{red} TODO:
% R * Btor towards axis - is constant? Attenberger check...
% }

\FloatBarrier
\section{Calculation of \texorpdfstring{$\nabla\alpha$}{grad-alpha} from VMEC equilibria}
%
But how do we get $\nabla\alpha$ or $\mathbf{\hat{e}}^\alpha$,
if not through the easily computable magnetic field direction and flux surface normal?
%
We can go via the angle definition, $\alpha = \theta - \iota\phi$ and get~\cite{plunk_collisionless_2014}:
\begin{align}
 \nabla\alpha = \nabla\theta - \iota\nabla\phi - \iota'\phi\nabla u_r \, ,
\end{align}
where $u_r$ is any (radial) coordinate and $\iota' = d\iota/du_r$.
%
Unfortunately, this is only true if the poloidal and toroidal angles are straight-fieldline-coordinates
(i.e. coordinates in which field lines are straight lines),
which is not the case for VMEC coordinates.
%
To transform to according coordinates, a stream function $\lambda(s,\theta_v, \phi)$ is necessary~\cite{stella_doc}:
\begin{align}
 \theta = \theta_v + \lambda(s,\theta_v, \phi) \, ,
\end{align}
where $\theta_v$ is the VMEC poloidal angle and
$\theta$ is a straight-fieldline-coordinate together with $\phi$ and $s$,
which are unchanged.
%
We then get $\alpha = \theta_v + \lambda - \iota\phi$ and accordingly~\cite{stella_doc}:
\begin{align}
    \nabla\alpha &= \nabla(\theta_v + \lambda - \iota\phi) \\
    &=   \bigg(\frac{\partial\lambda}{\partial s}
       - \phi\frac{d\iota}{ds}\bigg)\nabla s
       + \bigg(1 + \frac{\partial\lambda}{\partial\theta_v}\bigg)\nabla\theta_v
       + \bigg(-\iota + \frac{\partial\lambda}{\partial\phi}\bigg)\nabla\phi \, .
\end{align}
%
Find the relevant part in stella diagnostic code
\href{https://github.com/stellaGK/stella/blob/master/POST_PROCESSING/stellapy/data/geometry/calculate_geometricQuantitiesVMEC.py#L738}{here}.

The stream function $\lambda$ is expressed as a Fourier series:
\begin{align}
 \lambda(s, \theta_v, \phi) = \sum_{m=0}^M\sum_{n=-N}^{N}\hat{\lambda}_{m,n}^{\sin}(s)\sin{(m\theta - nN_\mathrm{fp}\phi)} \, ,
\end{align}
where $\hat{\lambda}_{m,n}^{\sin}(s)$ are the Fourier coefficients given in the VMEC output on a $(n_s, M\times 2N)$ grid.
%
From this, we get analytical expressions for the partial derivatives with respect to $\theta_v$ and $\phi$,
which can directly be calculated from the Fourier coefficients~\cite{stella_doc}:
\begin{align}
 \frac{\partial\lambda}{\partial\theta_v}(s, \theta_v, \phi) &=
   \sum_{m=0}^M\sum_{n=-N}^{N}
     m\hat{\lambda}_{m,n}^{\sin}(s)\cos{(m\theta - nN_\mathrm{fp}\phi)} \\
 \frac{\partial\lambda}{\partial\phi}(s, \theta_v, \phi) &=
   \sum_{m=0}^M\sum_{n=-N}^{N}
     -nN_\mathrm{fp}\hat{\lambda}_{m,n}^{\sin}(s)\cos{(m\theta - nN_\mathrm{fp}\phi)} \, .
\end{align}
%
The derivative with respect to $s$ can be pulled into the sum,
but needs be to calculated numerically based on the available $n_s$ grid of Fourier coefficients:
\begin{align}
 \frac{\partial\lambda}{\partial s}(s, \theta_v, \phi) &=
   \sum_{m=0}^M\sum_{n=-N}^{N}
     \bigg[ \frac{\partial}{\partial s}\hat{\lambda}_{m,n}^{\sin}(s) \bigg] \sin{(m\theta - nN_\mathrm{fp}\phi)} \\
 \frac{\partial}{\partial s}\hat{\lambda}_{m,n}^{\sin}(s) &\approx
   \frac{\hat{\lambda}_{m,n}^{\sin}[i_s +1] - \hat{\lambda}_{m,n}^{\sin}[i_s]}{\Delta s} \, .
\end{align}
